% Shared typography, theorem environments, notation, and navigation.
% All statement types share the theorem counter, reset at each chapter.
% Boxed definitions use that same counter; no independent defn counter exists.
\usepackage{fontspec}
\IfFontExistsTF{Times New Roman}{\setmainfont{Times New Roman}}{\setmainfont{Tinos}}
\IfFontExistsTF{Consolas}{\setmonofont{Consolas}[Scale=0.80]}{\setmonofont{DejaVu Sans Mono}[Scale=0.80]}

\usepackage{amsmath}
\usepackage{amssymb}
\usepackage{amsthm}
\usepackage{bm}
\usepackage{mathrsfs}
\usepackage{mathtools}

\usepackage{geometry}
\usepackage{setspace}
\usepackage{graphicx}
\usepackage{caption}
\usepackage{subcaption}
\usepackage{placeins}
\usepackage{booktabs}
\usepackage{enumitem}
\usepackage[bottom]{footmisc}
\usepackage{fancyhdr}

\usepackage[dvipsnames, svgnames]{xcolor}
\usepackage{pifont}
\usepackage{tcolorbox}
\tcbuselibrary{skins,breakable}

\usepackage{algorithm}
\usepackage[noend]{algpseudocode}
\usepackage{etoolbox}

\usepackage{ifthen}
\usepackage{hyperref}
\usepackage[backend=biber,style=alphabetic,sorting=nyt]{biblatex}
\addbibresource{ref.bib}
\AtBeginBibliography{\sloppy}
\usepackage{url}
\usepackage{float}
\usepackage{appendix}
\usepackage{animate}
\usepackage{tikz}

\definecolor{greenshade}{rgb}{0.9, 1, 0.9}
\definecolor{green}{rgb}{0.2, 0.6, 0.2}
\definecolor{axmBack}{rgb}{0.95, 0.95, 1}
\definecolor{axmFrame}{rgb}{0.5, 0.5, 1}
\definecolor{axmTitle}{rgb}{0.2, 0.2, 0.6}
\definecolor{lightblue}{rgb}{0.8, 0.9, 1}
\definecolor{lightgray}{rgb}{0.95, 0.95, 0.95}
\definecolor{lightyellow}{rgb}{1, 0.98, 0.9}
\definecolor{lightcyan}{rgb}{0.9, 1, 1}
\definecolor{SkyBlue}{rgb}{0.2,0.6,1}
\definecolor{lightcyanq}{RGB}{220, 245, 240}

\newtheoremstyle{mythm}
  {\topsep}{\topsep}{\itshape}{}{\bfseries}{.}{.5em}{}
\theoremstyle{mythm}
\newtheorem{theorem}{Theorem}[chapter]
\newtheorem{lemma}[theorem]{Lemma}
\newtheorem{corollary}[theorem]{Corollary}
\newtheorem{proposition}[theorem]{Proposition}
\newtheorem{axiom}[theorem]{Axiom}

\newtheoremstyle{mydef}
  {\topsep}{\topsep}{\upshape}{}{\bfseries}{.}{.5em}{}
\theoremstyle{mydef}

\newtheorem{example}[theorem]{Example}
\newtheorem{remark}[theorem]{Remark}

\renewenvironment{proof}{\indent\textcolor{SkyBlue}{\textbf{Proof.}}\;}{\qed\par}
\newenvironment{solution}{\indent\textcolor{SkyBlue}{\textbf{Solution.}}\;}{\qed\par}

\newtcolorbox[use counter=theorem, number within=chapter]{definition}[1][]{
  enhanced, breakable,
  colback=greenshade, colframe=green, arc=3mm,
  title={\textbf{Definition \thetcbcounter\ifstrempty{#1}{}{ (#1)}}},
  coltitle=black, fonttitle=\bfseries
}

\newtcolorbox[auto counter, number within=chapter]{axiombox}[1][]{
  enhanced, breakable,
  colback=axmBack, colframe=axmFrame, arc=3mm,
  title={\textbf{Axiom \thetcbcounter\ #1}},
  coltitle=axmTitle, fonttitle=\bfseries,
}

\newtcolorbox{Notice}[1][Notice]{
  colback=lightblue, colframe=black, arc=3mm,
  title={\textbf{#1}}, coltitle=white, fonttitle=\bfseries
}

\newtcolorbox{Proof}{
  colback=lightgray, colframe=black, arc=3mm,
  title={\textbf{Proof}}, coltitle=black, fonttitle=\bfseries
}

\newtcolorbox{Derivation}[1][]{
  enhanced, breakable,
  colback=lightyellow, colframe=black, arc=3mm,
  title={\textbf{Derivation}}, coltitle=white, fonttitle=\bfseries,
  boxsep=2mm, #1
}

\newtcolorbox{examplebox}[1][Example]{
  colback=lightcyan, colframe=black, arc=3mm,
  title={\textbf{#1}}, coltitle=black, fonttitle=\bfseries
}

\newtcolorbox{Question}[1][]{
  enhanced, breakable, colback=lightcyanq, colframe=black, arc=3mm,
  title={\textbf{Question}}, coltitle=white, fonttitle=\bfseries, boxsep=2mm, #1
}

\makeatletter
\newenvironment{breakablealgorithm}
  {%
    \begin{center}
      \refstepcounter{algorithm}%
      \hrule height.8pt depth0pt \kern4pt%
      \parskip 0pt
      \renewcommand{\caption}[2][\relax]{%
        {\raggedright\textbf{\fname@algorithm~\thealgorithm: ##2}\par}%
        \ifx\relax##1\relax
          \addcontentsline{loa}{algorithm}{\protect\numberline{\thealgorithm}##2}%
        \else
          \addcontentsline{loa}{algorithm}{\protect\numberline{\thealgorithm}##1}%
        \fi
        \kern2pt\hrule\kern2pt
     }
  }
  {%
     \kern2pt\hrule\relax%
   \end{center}
  }
\makeatother

\providecommand{\listofalgorithms}{\listof{algorithm}{List of Algorithms}}

% English manuscript additions: cross-referenced source annotations.
\usepackage{longtable}
\usepackage{array}
\geometry{left=29mm,right=27mm,top=28mm,bottom=28mm}
\setstretch{1.13}
% Show the thematic groups in the contents; local derivation headings stay in the text.
\setcounter{tocdepth}{2}
\setlength{\emergencystretch}{3em}
\hypersetup{hypertexnames=false,colorlinks=true,linkcolor=MidnightBlue,urlcolor=RoyalBlue,
  pdftitle={Probability Dynamics Lab: Stochastic Dynamics, Coarse-Graining, and Learned Probability Flows},
  pdfauthor={Research working notes}}
\theoremstyle{mydef}
\newtheorem{exercise}[theorem]{Exercise}
\newtheorem{assumption}[theorem]{Assumption}
\newcommand{\R}{\mathbb{R}}
\newcommand{\E}{\mathbb{E}}
\newcommand{\Pp}{\mathbb{P}}
\newcommand{\Law}{\operatorname{Law}}
\newcommand{\KL}{D_{\mathrm{KL}}}
\newcommand{\dd}{\,\mathrm{d}}
\newcommand{\supp}{\operatorname{supp}}
\newcommand{\sourcecredit}[2]{\par\smallskip\noindent{\footnotesize\itshape\textcolor{MidnightBlue}{Original exercise:} \emph{Student Assessment Exercise}\,\cite{assessment}, Project~#1, item~#2 (adapted).}\par\nobreak\smallskip\nobreak}
\newcommand{\sourcelink}[3]{\hyperref[#3]{Project #1 -- #2} & \hyperref[#3]{\pageref*{#3}}\\}
\newcommand{\sourceindexrow}[4]{\hyperref[#3]{\textcolor{MidnightBlue}{Project~#1 -- #2}} & \hyperref[#3]{\textcolor{MidnightBlue}{\pageref*{#3}}} & #4\\}
\pagestyle{fancy}
\fancyhf{}
\fancyhead[L]{\footnotesize Probability Dynamics Lab}
\fancyhead[R]{\footnotesize \nouppercase{\rightmark}}
\fancyfoot[C]{\thepage}
\renewcommand{\headrulewidth}{0.3pt}


% Display mathematics: keep text-size symbols and allow gentle page breaks.
\setlength{\jot}{5pt}
\allowdisplaybreaks[1]
\AtBeginDocument{%
  \setlength{\abovedisplayskip}{10pt plus 3pt minus 3pt}%
  \setlength{\belowdisplayskip}{10pt plus 3pt minus 3pt}%
  \setlength{\abovedisplayshortskip}{4pt plus 2pt}%
  \setlength{\belowdisplayshortskip}{8pt plus 2pt minus 2pt}%
}
